\section{主要工作总结}

本文围绕 LeRobot 在线控制中的``观测--推理--执行''闭环，研究了树莓派 5
这类无 GPU ARM 边缘平台上机器人基础软件栈的系统级性能问题。与只统计端到端
FPS 或只从模型、硬件通信单点优化入手的做法不同，本文将问题拆解到 AI 推理框架层、Python 运行时层、Linux
内核层与硬件 I/O 层，建立了面向机器人在线控制负载的四层性能剖析方法。
该方法不是依赖单一 profiling 工具给出结论，而是把 Python 阶段计时、
系统调用追踪、内核阻塞拆分和硬件通信路径分析相互对照，
从而为后续优化判断提供依据。

在实验环境方面，本文以 SO-101 六自由度机械臂、双摄像头输入和 ACT 策略模型为对象，
构建了可复现实验平台，并围绕核心控制循环记录各阶段耗时。
在测量工具层面，本文比较了 Python 阶段计时、strace
与 ftrace 的适用边界。工具链校准结果表明，
strace 会显著放大上下文切换并改变阶段时间结构，因此不适合承担定量时序分析；
ftrace 延迟导出方案对主循环扰动相对可控，更适合作为内核态观测工具。
在此基础上，本文进一步通过自定义内核插桩拆分了 pselect6 和
futex 路径中的阻塞时间与 CPU 处理时间，使用户态日志和内核态事件能够互相印证。

性能剖析结果显示，当前系统的主要瓶颈并不在单个系统调用或舵机写入函数本身，
而在 ARM CPU 上执行 ACT 模型前向推理的绝对耗时。观测阶段中，两路摄像头已经通过
OpenCV 后台线程完成异步预取，主线程读取缓存帧的开销较小；执行阶段中，
舵机同步读取与同步写入的毫秒级耗时主要来自
1 Mbps 半双工总线和 SCServo 协议约束，而不是 Linux 系统调用路径。
相比之下，单次完整 ACT 推理耗时约为秒级，其中 ResNet18 视觉特征提取和
Transformer Encoder 占据主要比例，是限制在线控制频率的核心因素。
这一结论说明，在当前硬件条件下，优化方向应优先指向推理运行时和模型前向计算耗时，
而不是盲目改写已受物理总线约束的硬件 I/O 路径。

在优化探索方面，本文首先分析了观测、推理与执行三类路径的异步化可行性。
由于舵机读写共享同一条半双工总线，简单后台线程化会退化为``持锁串行 + 线程切换''，
难以形成真正并行；相机路径已经具备后台预取机制，进一步收益有限。
因此，本文将优化重点放在推理异步化上，推导了异步流水线稳定无 stall 的近似条件
\(T_{\text{infer}}F \leq \min(hN,\;N/2)\)，并进一步建立了
FPS 随异步推理触发阈值在临界点前上升、临界点后饱和到目标帧率的分段模型。
实验结果表明，在 chunk\_size 为 100、目标 30 FPS 的设置下，
当前单轮扫描实验中表现最优的异步推理触发阈值0.30可取得 27.42 FPS 的有效控制频率，相对同步 baseline
提升 70.0\%，且稳态阶段几乎无 stall，说明异步调度能够显著回收动作执行期间的空窗。
episode 平均 FPS 尚未达到 30 的主要原因是启动空队列等待和树莓派 5 单机部署下的 CPU 资源争抢，
后续仍需从机器人基础软件栈中的推理运行时、操作系统调度和跨层诊断工具继续优化。
